% ==================================================================
%  RoundCount.tex  --  part of Section 3.1 (client-side preprocessing)
%  Kept in its own file because the rule is still being refined;
%  the rest of the paper refers to it only through \ref{sec:rounds}.
% ==================================================================
\subsubsection{Round Count}
\label{sec:rounds}

The last parameter is $R$, the number of rounds in which the server
spreads labels among the core points (Section~\ref{sec:lp}). Label
propagation has two parts, following the definition of DBSCAN. A cluster
grows only through its core points: a chain of neighboring core points
carries the cluster from one core point to the next, while a border
point joins the cluster of a core neighbor but extends it no further. If
a border point passed its label on, two clusters that share a border
point would merge, which DBSCAN does not allow. The server therefore
first spreads labels among the core points, round after round, and then
gives every non-core point the label of its core neighbors in a single
pass. Only the first part depends on how far a label has to travel. Each
round moves a label at least one hop along a chain of neighboring core
points, so $R$ rounds suffice once $R$ reaches the number of hops between
the two most distant core points of a cluster.

The client chooses $R$. Prior clustering over encrypted or secret-shared
data likewise leaves the number of iterations to the data owner. Cheon et
al.~\cite{cheon2019towards} choose the iteration counts separately for
each dataset, Yang et al.~\cite{yang2024optimized} take the number of
K-means iterations as a constant supplied by the data owner, and Bozdemir
et al.~\cite{bozdemir2021privacy} fix the number of expansion iterations
to four, which sufficed on their test data. None of them, however, gives
the client a way to choose that value in advance; Bozdemir et al.\ can
instead test for convergence during the run, at the cost of one bit of
leakage per test. The exact number of hops cannot be computed without the
neighbor relation and the core points, that is, without running DBSCAN on
the plaintext, which would leave nothing for the server to do. We
therefore give the client a naive estimate from a coarse grid.
It uses only cell indices and cell counts, never compares a distance with
$\varepsilon_n$, and gives the client a starting value when it chooses
$R$.

Let $s = \varepsilon_n / 2$ be the side of the grid, and assign each
point $\mathbf{x}$ to the cell $\lfloor \mathbf{x} / s \rfloor$,
coordinate by coordinate. Two points at distance at most
$\varepsilon_n = 2s$ differ by at most $2s$ along every axis, so their
cell indices differ by at most two along every axis. We call two cells
\emph{adjacent} when their indices differ by at most two along every
axis. The $5^d - 1$ cells adjacent to a cell then contain every possible
neighbor of its points. The client proceeds as follows.

\begin{enumerate}
\item \emph{Filter.} Drop every occupied cell whose own points and those
  of its adjacent cells number fewer than $\mathit{MinPts}$. Such a cell
  cannot contain a core point, since that count is an upper bound on the
  size of any $\varepsilon_n$-neighborhood inside it.
\item \emph{Group.} Connect every two remaining cells that are adjacent,
  and take the connected components as groups. The core points of every
  cluster lie inside a single group: two neighboring core points fall in
  adjacent cells, neither of which the filter drops, so a chain of
  neighboring core points never leaves its group.
\item \emph{Measure.} Give each edge of a group the distance between the
  centers of its two cells, and compute the geodesic diameter
  $\Lambda_g$, the longest shortest path inside the group.
\item \emph{Round up.} Take
  \begin{align*}
    R &= \max\Bigl(1, \; \max_{g}
      \Bigl\lceil \frac{\Lambda_g + s\sqrt{d}}{\varepsilon_n}
      \Bigr\rceil \Bigr) \\
      &= \max\Bigl(1, \; \max_{g}
      \Bigl\lceil \frac{\Lambda_g}{\varepsilon_n}
        + \frac{\sqrt{d}}{2} \Bigr\rceil \Bigr),
  \end{align*}
  where the term $s\sqrt{d}$ adds half a cell diagonal at each end of
  the path, because a point sits anywhere in its cell rather than at the
  center.
\end{enumerate}

The rule measures the length of the path a label has to travel, not the
straight-line extent of a cluster. The two differ on curved clusters
such as rings and spirals, where two core points can be close in a
straight line but far apart along the cluster. The rule has the
following property. For two core points $p$ and $q$ of a cluster in
group $g$, let $\mathrm{geo}(p,q)$ be the length of the shortest path
that leads from $p$ to the center of its cell, through the cell graph of
$g$ to the center of the cell of $q$, and on to $q$. The two end
segments are at most half a cell diagonal, $s\sqrt{d}/2$, each, and the
middle part is at most $\Lambda_g$, so
\[
  \Bigl\lceil \frac{\mathrm{geo}(p,q)}{\varepsilon_n} \Bigr\rceil
  \;\le\; \Bigl\lceil \frac{\Lambda_g + s\sqrt{d}}{\varepsilon_n}
  \Bigr\rceil \;\le\; R .
\]
$R$ is thus at least the number of steps of length $\varepsilon_n$ that
join any two core points of a cluster along a path through its group,
and $R$ rounds carry a label across every cluster whenever each round
moves it at least $\varepsilon_n$ along such a path.

A cell side of $\varepsilon_n/2$, rather than the
$\varepsilon_n/\sqrt{2}$ used in~\cite{fu2024ppa}, keeps the grid
binary-nested inside the $\varepsilon_n$ grid and makes adjacency a
plain two-cell offset along each axis. Let $C \le n$ be the number of
occupied cells and $T \le C$ the number of single-source shortest-path
runs used to find the diameters $\Lambda_g$. For a fixed dimension $d$,
so that every cell has the constant number $5^d - 1$ of adjacent cells,
the estimate takes $O(n + T\,C \log C)$ time, which is
$O(n + C^2 \log C)$ in the worst case.